% Standalone problem document, generated from the adjacent README.md.
% Compile with XeLaTeX. Mathematical content is maintained in Markdown.
\documentclass[11pt,a4paper]{article}
\usepackage[fontsize=10.5pt]{scrextend}
\usepackage[margin=22mm,headheight=15pt,headsep=9mm,footskip=11mm]{geometry}
\usepackage{fontspec}
\setmainfont{texgyrepagella}[Extension=.otf,UprightFont=*-regular,BoldFont=*-bold,ItalicFont=*-italic,BoldItalicFont=*-bolditalic]
\setsansfont{texgyreheros}[Extension=.otf,UprightFont=*-regular,BoldFont=*-bold,ItalicFont=*-italic,BoldItalicFont=*-bolditalic]
\setmonofont{texgyrecursor}[Extension=.otf,UprightFont=*-regular,BoldFont=*-bold,ItalicFont=*-italic,BoldItalicFont=*-bolditalic,Scale=MatchLowercase]
\usepackage{amsmath,amssymb}
\usepackage{unicode-math}
\setmathfont{texgyrepagella-math.otf}
\usepackage{xcolor,microtype,enumitem,needspace,fancyhdr}
\usepackage{bookmark}
\definecolor{ink}{HTML}{17344B}
\definecolor{accent}{HTML}{197D80}
\definecolor{muted}{HTML}{596773}
\hypersetup{colorlinks=true,linkcolor=accent,urlcolor=accent,pdftitle={RA-12 - Relative-error
threshold for extremal Gaussian trace
bounds},pdfauthor={Open Problems in Numerical Linear Algebra}}
\urlstyle{same}
\setlength{\parindent}{0pt}
\setlength{\parskip}{5pt plus 1pt minus 1pt}
\linespread{1.04}
\setlength{\emergencystretch}{3em}
\widowpenalty=10000
\clubpenalty=10000
\displaywidowpenalty=10000
\allowdisplaybreaks[1]
\setlist{nosep,leftmargin=1.5em}
\providecommand{\tightlist}{\setlength{\itemsep}{0pt}\setlength{\parskip}{0pt}}
\providecommand{\pandocbounded}[1]{#1}
\pagestyle{fancy}
\fancyhf{}
\fancyhead[L]{\sffamily\footnotesize\color{muted}OPEN PROBLEMS IN NUMERICAL LINEAR ALGEBRA}
\fancyhead[R]{\sffamily\bfseries\footnotesize\color{accent}RA-12}
\fancyfoot[L]{\parbox[b]{0.9\textwidth}{\sffamily\fontsize{8}{10}\selectfont\color{muted}Editorial ratings; a bounded search cannot certify that a problem remains unsolved.\\Literature check: 2026-09-11}}
\fancyfoot[R]{\sffamily\footnotesize\color{muted}\thepage}
\renewcommand{\headrulewidth}{0.3pt}
\renewcommand{\headrule}{\hbox to\headwidth{\color{accent}\leaders\hrule height \headrulewidth\hfill}}
\makeatletter
\renewcommand{\section}{\Needspace{5\baselineskip}\@startsection{section}{1}{0pt}{12pt plus 2pt minus 2pt}{4pt}{\sffamily\bfseries\large\color{ink}}}
\renewcommand{\subsection}{\Needspace{4\baselineskip}\@startsection{subsection}{2}{0pt}{9pt plus 1pt minus 1pt}{3pt}{\sffamily\bfseries\normalsize\color{ink}}}
\makeatother
\setcounter{secnumdepth}{0}
\begin{document}
{\sffamily\small\color{muted}Randomized and low-rank approximation\par}
\vspace{3pt}
{\raggedright\hyphenpenalty=10000\exhyphenpenalty=10000\sffamily\bfseries\fontsize{21}{24}\selectfont\color{ink}Relative-error
threshold for extremal Gaussian trace bounds\par}
\vspace{7pt}
{\sffamily\small\textbf{Difficulty:} challenging\quad\textbf{Importance:} interesting
to the community\par}
{\sffamily\small\bfseries\color{accent}Status: Open\par}
\vspace{4pt}
{\color{accent}\hrule height 0.8pt}
\vspace{6pt}
\textbf{Rating rationale:} Challenging because an explicit threshold
must control extremal tails uniformly over spectra; community impact is
sharper distribution-level trace-estimation confidence bounds.\\
\textbf{Source:} Hallman, Conjecture 3 together with Theorem 6.

\subsection{Related auxiliary counterexamples ---
2026-09-11}\label{related-auxiliary-counterexamples-2026-09-11}

\textbf{Author:} Matthew J. Colbrook, Department of Applied Mathematics
and Theoretical Physics, University of Cambridge. \textbf{Independent
Codex-agent review: PASS for the limited result below.}

The submitted Gamma-density examples refute the upper-mode assertion in
Hallman Conjecture 1 and the upper-inflection assertion in Conjecture 2,
with legally distinct augmentation indices. These are counterexamples to
auxiliary assertions.

\textbf{Remaining question:} The complete relative Gaussian trace-tail
probability chain in this entry is neither proved nor refuted. No
revised sharp tail threshold is established; status remains Open.

\textbf{Primary reference:}
\href{https://github.com/ajt60gaibb/OpenProblemsInNLA/blob/main/references/colbrook-transfer-2026-09-11/manuscripts/04_gamma_auxiliary_counterexamples.pdf}{complete
authored PDF},
\href{https://github.com/ajt60gaibb/OpenProblemsInNLA/blob/main/references/colbrook-transfer-2026-09-11/manuscripts/04_gamma_auxiliary_counterexamples.tex}{standalone
TeX}, \textbf{Proposition 2.1 (with Proposition 3.1 for related
evidence)}.
\href{https://github.com/ajt60gaibb/OpenProblemsInNLA/blob/main/references/colbrook-transfer-2026-09-11/verification/reviews/gamma-auxiliary-review.md}{Independent
proof review} ·
\href{https://github.com/ajt60gaibb/OpenProblemsInNLA/blob/main/references/colbrook-transfer-2026-09-11/README.md}{Authorship
and submission record}. Verification is independent agent review, not
external human peer review or formal certification.

\subsection{Problem statement}\label{problem-statement}

For a real symmetric \(d\times d\) matrix \(D\) and integer \(m\ge1\),
define the Gaussian trace estimator

\[
T_m(D)=\frac1m\sum_{j=1}^{m}z_j^TDz_j,
\qquad z_1,\ldots,z_m\overset{\mathrm{iid}}{\sim}N(0,I_d).
\]

Let \(A\ne0\) be any real symmetric positive semidefinite \(n\times n\)
matrix, with \(n\ge1\), and set

\[
\mu=\frac{\operatorname{tr}(A)}{\|A\|_2},\qquad
B_\mu=\frac1\mu\operatorname{diag}
\left(I_{\lfloor\mu\rfloor},\,\mu-\lfloor\mu\rfloor\right).
\]

Here \(\|\cdot\|_2\) is the spectral norm, \(\mu\ge1\), and a zero final
diagonal entry may be retained. Let \(X\) have the Gamma distribution
with shape and rate both \(m\mu/2\). The shape/rate convention means
that \(\operatorname{Gamma}(\alpha,\beta)\) has density
\(\beta^\alpha x^{\alpha-1}e^{-\beta x}/\Gamma(\alpha)\) for \(x>0\).

\textbf{Conjecture.} For every such \(A\), every integer \(m\ge1\), and
every \(\varepsilon\ge2/(m\mu)\), the complete comparison chain holds:\\
\[
\begin{aligned}
\Pr\!\left(
 |T_m(A)-\operatorname{tr}(A)|
 \ge\varepsilon\operatorname{tr}(A)
\right)
&\le \Pr\!\left(|T_m(B_\mu)-1|\ge\varepsilon\right)\\
&\le \Pr\!\left(|X-1|\ge\varepsilon\right).
\end{aligned}
\]

Each estimator uses Gaussian vectors of its own matrix dimension; only
their distributions are compared.

\subsection{Why it matters in numerical linear
algebra}\label{why-it-matters-in-numerical-linear-algebra}

This would specify the tolerance range on which effective rank yields
extremal, distribution-level confidence bounds for Gaussian trace
estimation. It also applies to Frobenius-norm estimation through
\(\|C\|_F^2=\operatorname{tr}(C^TC)\).

\newpage
\subsection{References}\label{references}

\begin{enumerate}
\def\labelenumi{\arabic{enumi}.}
\tightlist
\item
  Eric Hallman,
  \href{https://arxiv.org/html/2411.15454v1\#S5}{\emph{Extremal bounds
  for Gaussian trace estimation}}, arXiv:2411.15454v1 (2024), §5,
  Theorem 6 and Conjecture 3.
\item
  Alice Cortinovis and Daniel Kressner,
  \href{https://doi.org/10.1007/s10208-021-09525-9}{\emph{On Randomized
  Trace Estimates for Indefinite Matrices with an Application to
  Determinants}}, FoCM 22 (2022), 875--903, Theorem 1.
\end{enumerate}

\subsection{Status check}\label{status-check}

Theorem 6 proves the comparisons beyond an unspecified threshold;
Conjecture 3 supplies the explicit threshold above. On 2026-09-08 the
\href{https://arxiv.org/abs/2411.15454}{arXiv record} still listed only
v1. Title, author, conjecture-number and 2025/2026 searches found no
resolution. The author's \href{https://arxiv.org/abs/2512.02316}{later
XTrace paper} studies different estimators. This is a bounded check.

\subsection{Audit --- 2026-09-10}\label{audit-2026-09-10}

Rechecked \href{https://arxiv.org/html/2411.15454v1}{Hallman, Theorem 6
and Conjecture 3}; the record still lists only v1. Author,
Gaussian-trace, and conjecture searches found no resolution. An
unspecified larger threshold proves a weaker result and does not verify
the displayed explicit threshold.
\end{document}
